\chapter{Smooth-gradient and affine realizations}
\label{chap:smooth-affine}

This chapter instantiates the abstract Basic.  The direct realization lives in the ambient
Euclidean space.  The intrinsic realization first removes null normal directions by passing to
the smallest closed affine support.  The repository selects the latter realization as its
canonical frontend; this is a design choice, not a hidden identification with a literature
Sobolev or Cheeger energy.

\section{Ambient smooth-gradient relaxation}

\begin{definition}[Smooth-gradient pre-energy]
  \label{def:smooth-gradient-preenergy}
  \lean{FunctionalInequalities.Energy.Euclidean.SmoothGradientTest,
        FunctionalInequalities.Energy.Euclidean.SmoothGradientTest.gradientEnergy,
        FunctionalInequalities.Energy.Euclidean.SmoothGradientTest.toL2,
        FunctionalInequalities.Energy.Euclidean.smoothGradientPreEnergy}
  \leanok
  For a Borel probability $\mu$ on a finite-dimensional real inner-product space, the test space
  consists of globally $C^\infty$ real functions with an $L^2(\mu)$ realization and integrable
  squared gradient.  The pre-energy is
  $\mathcal E_0(u)=\int\norm{\nabla u}^2\dd\mu$.
\end{definition}

\begin{theorem}[Quadratic dense core]
  \label{thm:smooth-gradient-quadratic-dense}
  \lean{FunctionalInequalities.Energy.Euclidean.smoothGradientPreEnergy_quadratic,
        FunctionalInequalities.Energy.Euclidean.smoothGradientPreEnergy_denseRealization}
  \leanok
  \uses{def:smooth-gradient-preenergy}
  The smooth-gradient pre-energy is quadratic, and its realization is dense in real $L^2(\mu)$.
\end{theorem}
\begin{proof}
  \leanok
  The gradient is linear and its squared norm satisfies the parallelogram law.  Density is the
  finite-dimensional Borel $L^2$ density theorem for compactly supported smooth functions.
\end{proof}

\begin{definition}[Ambient relaxed gradient energy]
  \label{def:smooth-gradient-relaxation}
  \lean{FunctionalInequalities.Energy.Euclidean.smoothGradientRelaxation}
  \leanok
  \uses{def:preenergy-relaxation,def:smooth-gradient-preenergy}
  The ambient smooth-gradient energy is the sequential $L^2(\mu)$ relaxation of the preceding
  pre-energy.
\end{definition}

\begin{theorem}[Ambient relaxation is closed quadratic]
  \label{thm:smooth-gradient-relaxation-closed}
  \lean{FunctionalInequalities.Energy.Euclidean.smoothGradientRelaxation_closedQuadratic}
  \leanok
  \uses{thm:quadratic-relaxation-closed,thm:smooth-gradient-quadratic-dense,def:smooth-gradient-relaxation}
  The ambient smooth-gradient relaxation is a densely defined closed quadratic energy on
  $L^2(\mu)$.
\end{theorem}
\begin{proof}
  \leanok
  Apply the abstract quadratic-relaxation theorem to the smooth quadratic dense core.
\end{proof}

\section{Affine support and intrinsic energy}

\begin{definition}[Affine support and intrinsic law]
  \label{def:affine-support-intrinsic-law}
  \lean{FunctionalInequalities.Probability.affineSupport,
        FunctionalInequalities.Probability.intrinsicMap,
        FunctionalInequalities.Probability.intrinsicMeasure,
        FunctionalInequalities.Probability.reconstruct}
  \leanok
  The affine support of $\mu$ is the smallest closed affine subspace containing almost every
  point.  Choosing its canonical direction space gives an intrinsic coordinate map, the pushed
  intrinsic law, and a reconstruction map back to the ambient space.
\end{definition}

\begin{theorem}[Intrinsic $L^2$ equivalence]
  \label{thm:intrinsic-l2-equivalence}
  \lean{FunctionalInequalities.Energy.Euclidean.toIntrinsicL2,
        FunctionalInequalities.Energy.Euclidean.toAmbientL2,
        FunctionalInequalities.Energy.Euclidean.toAmbientL2_toIntrinsicL2,
        FunctionalInequalities.Energy.Euclidean.toIntrinsicL2_toAmbientL2,
        FunctionalInequalities.Energy.Euclidean.intrinsicL2Equiv}
  \leanok
  \uses{def:affine-support-intrinsic-law,def:l2-mean-variance}
  Pullback and reconstruction define inverse real linear isometries between ambient
  $L^2(\mu)$ and intrinsic $L^2$ of the intrinsic law.
\end{theorem}
\begin{proof}
  \leanok
  Both coordinate maps are measure preserving, and reconstruction is almost everywhere inverse
  to the intrinsic map on the affine support.
\end{proof}

\begin{definition}[Intrinsic-affine relaxed energy]
  \label{def:intrinsic-affine-energy}
  \lean{FunctionalInequalities.Energy.Euclidean.intrinsicAffineSmoothGradientRelaxation}
  \leanok
  \uses{def:affine-support-intrinsic-law,thm:intrinsic-l2-equivalence,def:smooth-gradient-relaxation}
  The intrinsic-affine energy of an ambient $L^2(\mu)$ class is the intrinsic law's
  smooth-gradient relaxation evaluated after the intrinsic $L^2$ isometry.
\end{definition}

\begin{theorem}[Intrinsic-affine energy is closed quadratic]
  \label{thm:intrinsic-affine-energy-closed}
  \lean{FunctionalInequalities.Energy.Euclidean.intrinsicAffineSmoothGradientRelaxation_closedQuadratic}
  \leanok
  \uses{thm:smooth-gradient-relaxation-closed,thm:intrinsic-l2-equivalence,def:intrinsic-affine-energy}
  The intrinsic-affine relaxed energy is a densely defined closed quadratic energy on ambient
  $L^2(\mu)$.
\end{theorem}
\begin{proof}
  \leanok
  Transport the intrinsic closed quadratic form along the real linear isometric equivalence.
\end{proof}

\begin{definition}[Selected canonical frontend]
  \label{def:selected-canonical-frontend}
  \lean{FunctionalInequalities.Energy.Euclidean.intrinsicAffineSmoothGradientRelaxation,
        FunctionalInequalities.Poincare.Euclidean.intrinsicAffineSmoothGradientRelaxationConstant,
        FunctionalInequalities.Poincare.Euclidean.IntrinsicAffineSmoothGradientRelaxationHolds}
  \leanok
  \uses{def:intrinsic-affine-energy,def:l2-energy-poincare-constant}
  The selected canonical energy and constant are
  \[
    \mathcal E^{\mathrm{can}}_\mu
      :=\mathcal E^{\mathrm{intr}}_\mu,
    \qquad
    \CPCan(\mu)
      :=\CE\bigl(\mu,\mathcal E^{\mathrm{can}}_\mu\bigr).
  \]
  The corresponding finite-constant predicate is the $L^2$ energy inequality for
  $\mathcal E^{\mathrm{can}}_\mu$.
\end{definition}

The word ``selected'' is part of the contract.  This definition does not assert
$\mathcal E^{\mathrm{can}}_\mu=2\Ch_\mu$, equality with the ambient presentation, or equality with
the locally Lipschitz literature constant.  Those are separate theorem nodes.

\begin{theorem}[Intrinsic-affine optimal constant law]
  \label{thm:intrinsic-affine-constant-order}
  \lean{FunctionalInequalities.Poincare.Euclidean.intrinsicAffineSmoothGradientRelaxationConstant,
        FunctionalInequalities.Poincare.Euclidean.IntrinsicAffineSmoothGradientRelaxationHolds,
        FunctionalInequalities.Poincare.Euclidean.intrinsicAffineSmoothGradientRelaxationHolds_iff_constant_le}
  \leanok
  \uses{def:selected-canonical-frontend,thm:l2-energy-holds-order}
  The intrinsic-affine Poincar\'e predicate holds at $C$ exactly when its optimal extended
  constant is at most $C$.
\end{theorem}
\begin{proof}
  \leanok
  This is the canonical $L^2$ order law instantiated with the intrinsic-affine energy.
\end{proof}
